\documentclass[11pt]{amsart}
\usepackage[T1]{fontenc}
\usepackage{lmodern}
\usepackage{amsmath,amssymb,amsthm,mathtools}
\usepackage{microtype,hyperref}
\newtheorem{theorem}{Theorem}
\newtheorem{lemma}[theorem]{Lemma}
\theoremstyle{definition}
\newtheorem{definition}[theorem]{Definition}
\title{A Synthetic Mathematical Note}
\author{A. Author}
\address{Department of Reproducible Typesetting}
\email{author@example.invalid}
\date{September 2001}

\begin{document}
\begin{abstract}
We record a compact theorem, alignment, matrix, and bibliography-free
mathematical workload.
\end{abstract}
\maketitle

\section{The statement}
\begin{definition}
A sequence $(a_n)$ is stable if $a_{n+1}=a_n$ for all sufficiently large $n$.
\end{definition}

\begin{lemma}\label{lem:finite}
Every eventually constant sequence has a finite set of changes.
\end{lemma}

\begin{proof}
Choose $N$ after which the sequence is constant. Only indices below $N$ can
contribute a change.
\end{proof}

\begin{theorem}
For $A=\begin{pmatrix}1&1\\0&1\end{pmatrix}$, one has
\[
  A^n=\begin{pmatrix}1&n\\0&1\end{pmatrix}.
\]
\end{theorem}

\begin{align}
 f(x+y) &= f(x)+f(y), \\
 \frac{d}{dx}\int_0^x f(t)\,dt &= f(x).
\end{align}
See Lemma~\ref{lem:finite}.
\end{document}
